\documentclass[11pt,a4paper]{article}

% --- Codificación e idioma ---
\usepackage[utf8]{inputenc}
\usepackage[T1]{fontenc}
\usepackage[spanish,es-tabla]{babel}
\usepackage{csquotes}

% --- Geometría y formato ---
\usepackage[margin=2.5cm]{geometry}
\usepackage{setspace}
\onehalfspacing

% --- Gráficos, tablas y código ---
\usepackage{graphicx}
\usepackage{booktabs}
\usepackage{float}
\usepackage{listings}
\usepackage{xcolor}
\usepackage{amsmath}
\usepackage{hyperref}

\hypersetup{
    colorlinks=true,
    linkcolor=black,
    citecolor=blue,
    urlcolor=blue
}

% --- Bibliografía: Estilo IEEE estricto con BibLaTeX (Norma 4 y Criterio P4) ---
\usepackage[backend=biber,style=ieee,sorting=none]{biblatex}
\addbibresource{referencias.bib}

% --- Datos del Documento (Completar según la guía de la actividad) ---
\title{\textbf{FORM-MC-01: Mapa conceptual colaborativo}\\[0.3cm]
\large Universidad Técnica Estatal de Quevedo\\
Facultad de Ciencias de la Computación -- Carrera de Software\\
Aplicaciones Distribuidas (ISR-701) -- Séptimo Nivel ``A''}

\author{[Nombre del Estudiante 1] \\ \texttt{[correo@uteq.edu.ec]}}
\date{\today}

\begin{document}

\maketitle

\begin{abstract}
% Escribir aquí un resumen propio una vez finalizado el trabajo.
[Espacio en blanco para el resumen ejecutivo de la actividad.]
\end{abstract}

\section{Introducción y Fundamentación Conceptual (D2)}
% Desarrollar aquí los conceptos con redacción propia y citas reales, por ejemplo: \cite{vansteen2023}.
[Espacio en blanco para desarrollar la fundamentación conceptual de la unidad.]

\section{Desarrollo Técnico y Evidencia (D3)}
% Documentar aquí la implementación, arquitectura, configuración, código en src/ y datos en data/.
[Espacio en blanco para documentar el desarrollo técnico y referenciar archivos del repositorio.]

\section{Análisis Crítico y Resultados (D4)}
% Interpretar mediciones, comparar alternativas (trade-offs), reconocer limitaciones.
[Espacio en blanco para las tablas de mediciones, gráficas de figs/ y análisis crítico.]

\section{Conclusiones}
% Conclusiones cuantificables derivadas directamente de las pruebas realizadas.
[Espacio en blanco para las conclusiones propias del trabajo.]

% --- SECCIÓN OBLIGATORIA SEGÚN SECCIÓN 6.1 y 7.5 DEL ENCUADRE ---
% Nota: La falta de esta declaración resta 0.3 puntos sobre 10.
\section*{Declaración de Uso de Inteligencia Artificial}
% Elegir Opción A o B según corresponda a tu trabajo real en esta actividad:

% OPCIÓN A (Si NO se utilizó ninguna herramienta de IA):
% El autor (o equipo) declara que en la elaboración de esta actividad no se utilizó ninguna herramienta de inteligencia artificial generativa ni de asistencia automatizada. Todo el contenido, código, datos y redacción son de autoría íntegramente propia.

% OPCIÓN B (Si se utilizó IA únicamente como apoyo permitido):
Para la realización de esta actividad se utilizó la herramienta [Nombre de la herramienta, ej. GitHub Copilot / ChatGPT / Gemini] (versión [Versión/Modelo]) únicamente como herramienta de apoyo para [especificar uso permitido: p. ej., revisión ortográfica / consulta de errores de compilación en LaTeX]. El resultado obtenido fue verificado mediante [explicar cómo se comprobó técnicamente, p. ej., compilación local en TeX Live y contraste con la documentación oficial]. Todo el análisis, diseño, código evaluable, mediciones y redacción técnica corresponden a la autoría propia del estudiante/equipo.

% --- Impresión de Referencias (Estilo IEEE) ---
\printbibliography[title={Referencias}]

\end{document}
